
\chapter{Análise Integrada de Sistemas Biológicos}
\label{cap:11}

\normalfont\rmfamily\upshape
A biologia tem, há mais de quatro décadas, mudado a unidade de observação. Da
gene isolado --- observado nas décadas 1960--1970, quando a tecnologia
permitia marcar e seguir uma proteína --- passamos para redes gênicas nos anos
1980--1990, depois para vias inteiras com o sequenciamento estagiado da
\textit{Drosophila} e do projeto genoma humano (1990--2003), e finalmente
para a topologia completa do transcriptoma, do proteoma, do metaboloma e do
fluxoma de uma célula ou tecido sob condição definida. Esse acúmulo de
camadas informacionais simultâneas define a chamada \emph{biologia de sistemas}
e impõe um problema novo: como extrair regularidades dinâmicas a partir de
dados fragmentários, ruidosos e de dimensionalidade colossal. A resposta
disciplinada a essa pergunta é o tema deste capítulo.

O capítulo tem quatro objetivos centrais. \textbf{Primeiro}, descrever a
estratégia de integração multi-ômica em suas variantes metodológicas: correlação,
redes, modelos baseados em restrições, aprendizado de máquina. \textbf{Segundo},
ilustrar esse ferramental em um sistema-modelo bem caracterizado --- o ciclo da
trealose em \textit{Saccharomyces cerevisiae} --- onde a integração de RNA-seq,
proteômica e metabolômica sob choque térmico permite validar experimentalmente
previsões do modelo. \textbf{Terceiro}, discutir a reconstrução de vias e
redes a partir de dados, com foco em análise de enriquecimento, GSEA, redes
multi-camada e detecção de módulos. \textbf{Quarto}, apresentar abordagens de
modelagem matemática integrada --- GEMs (modelos em escala genômica), FBA,
modelos cinéticos e híbridos --- aplicadas à interpretação quantitativa de
conjuntos multi-ômicos.

\section{Introdução à Biologia de Sistemas}
\label{sec:11-introducao}

\begin{definicaoenv}[Biologia de Sistemas]
Biologia de sistemas é o estudo de sistemas biológicos mediante integração
quantitativa de dados multi-ômicos para compreender propriedades emergentes
e o comportamento do sistema como um todo. Sua hipótese central é que
componentes moleculares são organizados em redes modulares, cujas
propriedades --- robustez, oscilação, biestabilidade --- não podem ser
inferidas pela simples inspeção das partes.
\end{definicaoenv}

A biologia de sistemas encontra sua identidade formal em contraste com a
genômica funcional clássica: enquanto esta procura relações \emph{localizadas}
(mutação $\to$ fenótipo, proteína $\to$ função), aquela procura relações
\emph{globais} --- entender como o sistema \emph{se organiza} para responder a
uma perturbação, e como esse padrão de organização pode ser medido,
quantificado e previsto. Para responder a essas perguntas, três ferramentas
se mostraram particularmente úteis: medições multi-camada em séries
temporais, integração computacional rigorosa, e modelos matemáticos que
oferecem uma hipótese explícita e falsificável sobre o sistema. A
Figura~\ref{fig:11-multiescala} ilustra esquematicamente o escopo da
integração multi-ômica: as cinco camadas ômicas --- genômica,
transcriptômica, proteômica, metabolômica e fenômica --- alimentam
modelos em quatro escalas biológicas --- molecular, celular, tecidual e
do organismo ---, em uma pirâmide de integração que conecta os dados
com a função sistêmica.

\begin{figure}[htbp]
\centering
\begin{tikzpicture}[scale=0.8, every node/.style={font=\small}]
    % Multi-scale integration
    \node[draw, circle, fill=roxodna!30, text width=2cm, align=center] at (0,3) {Molecular};
    \node[draw, circle, fill=laranjacel!30, text width=2cm, align=center] at (3,3) {Celular};
    \node[draw, circle, fill=verdeacido!30, text width=2cm, align=center] at (6,3) {Tecido};
    \node[draw, circle, fill=azulclaro!30, text width=2cm, align=center] at (9,3) {Organismo};

    % Arrows showing integration
    \draw[->, ultra thick] (1,3) -- (2,3);
    \draw[->, ultra thick] (4,3) -- (5,3);
    \draw[->, ultra thick] (7,3) -- (8,3);

    % Bottom layer: omics
    \node[font=\tiny] at (0,1.5) {Genômica};
    \node[font=\tiny] at (2,1.5) {Transcriptômica};
    \node[font=\tiny] at (4,1.5) {Proteômica};
    \node[font=\tiny] at (6,1.5) {Metabolômica};
    \node[font=\tiny] at (8,1.5) {Fenômica};

    \draw[<->, thick, azulescuro] (0,2) -- (8,2);
    \node[above, font=\footnotesize] at (4,2) {Integração Multi-Escala};
\end{tikzpicture}
\caption{Esquema conceitual de integração multiômica em múltiplas escalas. As cinco camadas de dados ômicos (genômica, transcriptômica, proteômica, metabolômica, fenômica --- embaixo) alimentam modelos em quatro níveis de organização biológica (molecular, celular, tecidual, organismo --- em cima). As setas entre níveis indicam que a integração permite transferir informação de uma escala para a outra --- por exemplo, inferir funções celulares a partir de dados moleculares, ou fenótipos organismais a partir de assinaturas teciduais.}
\label{fig:11-multiescala}
\end{figure}

\subsection{Estratégia de Integração Multi-Ômica}
\label{sub:image11-estrategia}

A tabela central da biologia de sistemas é a chamada \emph{camadaômica} ---
uma matriz cujas linhas são \emph{features} (genes, transcritos, proteínas,
metabólitos) e cujas colunas são amostras sob diferentes condições. Quando
várias camadas são medidas na mesma série de amostras, surge a possibilidade
de integrar os dados para extrair sinais que nenhuma camada isolada mostraria.

Cinco camadas são as mais estudadas atualmente: \textbf{genômica} (variantes,
número de cópias, mutações), \textbf{transcriptômica} (RNA-seq, microarrays,
RNA-seq de célula única), \textbf{proteômica} (espectrometria de massas,
Western blot quantitativo), \textbf{metabolômica} (LC/MS ou NMR de metabólitos
polares e lipídios) e \textbf{fluxômica} ($^{13}$C-MFA, dados de distribuição
de massa de isótopos marcados). Cada camada tem dinâmica característica:
genômica é estática em escala tecidual; transcriptômica, com meia-vida da
ordem de horas; proteômica, com meia-vida de horas a dias; metabolômica e
fluxômica, com tempos de segundos a minutos.

A integração efetiva dessas camadas impõe cinco desafios computacionais e
conceituais. \textbf{Primeiro}, escalas dispares: RNAs são contados em
escala absoluta ou TPM, proteínas em intensidade iBAQ, metabólitos em
nanomolar ou milimolar. \textbf{Segundo}, dimensionalidade: o número de
features varia de $10^4$ (RNAs) a $10^2$ (metabólitos). \textbf{Terceiro},
sparsidade: muitos zeros legítimos (genes não-expressos, proteínas abaixo do
limite de detecção). \textbf{Quarto}, ruído técnico: preparação de amostra,
lote de reagente, profundidade de sequenciamento. \textbf{Quinto}, ruído
biológico: variabilidade inter-individual e inter-tempo. Resolver esses cinco
desafios é o trabalho prático do bioinformata de sistemas. A
Figura~\ref{fig:11-camadas-omicas} sintetiza as camadas ômicas em uma
pilha vertical análoga ao dogma central: do genoma (estático) ao
transcriptoma (resposta rápida), ao proteoma (função executora) e ao
metaboloma (fenótipo mais próximo).

\begin{figure}[htbp]
\centering
\begin{tikzpicture}[scale=0.85]
    % Central dogma with omics layers
    \node[draw, rectangle, fill=roxodna!30, minimum width=2.5cm, minimum height=1cm] at (0,4) {Genômica};
    \node[draw, rectangle, fill=laranjacel!30, minimum width=2.5cm, minimum height=1cm] at (0,2.5) {Transcriptômica};
    \node[draw, rectangle, fill=verdeacido!30, minimum width=2.5cm, minimum height=1cm] at (0,1) {Proteômica};
    \node[draw, rectangle, fill=azulclaro!30, minimum width=2.5cm, minimum height=1cm] at (0,-0.5) {Metabolômica};

    % Arrows showing flow
    \draw[->, ultra thick] (0,3.5) -- (0,3.0) node[right, midway, font=\tiny] {transcrição};
    \draw[->, ultra thick] (0,2.0) -- (0,1.5) node[right, midway, font=\tiny] {tradução};
    \draw[->, ultra thick] (0,0.5) -- (0,0) node[right, midway, font=\tiny] {catálise};

    % Characteristics on right
    \node[right, font=\tiny, text width=3.5cm] at (2,4) {DNA: estático, variações\\SNPs, CNVs, mutações};
    \node[right, font=\tiny, text width=3.5cm] at (2,2.5) {mRNA: dinâmico, 10$^3$-10$^7$ cópias\\responde rapidamente};
    \node[right, font=\tiny, text width=3.5cm] at (2,1) {Proteínas: funcional, modificações\\10$^4$-10$^6$ moléculas};
    \node[right, font=\tiny, text width=3.5cm] at (2,-0.5) {Metabólitos: fenotípico\\10$^{-9}$-10$^{-3}$ M};
\end{tikzpicture}
\caption{Camadas ômicas em analogia ao dogma central. Genômica (estática, escala tecidual), transcriptômica (dinâmica, 10$^3$--10$^7$ cópias), proteômica (10$^4$--10$^6$ moléculas, com modificações pós-traducionais) e metabolômica (10$^{-9}$--10$^{-3}$ M) são lidas em série: cada camada adiciona informação complementar sobre o estado do sistema. A seta de catálise --- proteína $\to$ metabólito --- é a única via que modifica o ambiente químico; as outras são informacionais.}
\label{fig:11-camadas-omicas}
\end{figure}

\section{Integração Multi-ômica}
\label{sec:11-multio-mica}

\subsection{Estratégias de Fusão: Early, Intermediate e Late}
\label{sub:image11-fusao}

A escolha metodológica mais decisiva em integração é a estratégia de \emph{fusão}.
Três famílias principais são distinguidas pela literatura:

\begin{itemize}
\item \textbf{Early fusion} (concatenação): todas as camadas são unificadas em
  uma matriz única, e o método de análise trata todas as features
  indistintamente. Vantagem: simplicidade de implementação. Desvantagem:
  nenhuma proteção contra a heterogeneidade entre camadas --- genes
  altamente variáveis podem dominar o sinal.
\item \textbf{Intermediate fusion} (intermediário): cada camada é reduzida
  independentemente a uma representação latente (PCA, autoencoder) e as
  representações são então combinadas em um novo modelo unificado. Aqui se
  situam os métodos de fatoração multi-ômica como \emph{MOFA}
  (Multi-Omics Factor Analysis), que identificam fatores latentes compartilhados.
\item \textbf{Late fusion} (consenso): cada camada é modelada separadamente até
  o fim, e os resultados são então combinados por voto, média ou regra
  bayesiana. Robusto a problemas específicos de uma camada; mais conservador.
\end{itemize}

A Figura~\ref{fig:11-fusao} compara as três estratégias graficamente:
na \emph{early fusion} todas as camadas são concatenadas em uma única
matriz integrada; na \emph{intermediate fusion} cada camada é reduzida
a representações latentes (módulos) que depois se combinam; na \emph{late
fusion} cada camada é modelada independentemente e os modelos produzem
um consenso por voto, média ou regra bayesiana.

\begin{figure}[htbp]
\centering
\begin{tikzpicture}[scale=0.7, every node/.style={font=\tiny}]
    % Early fusion
    \node[font=\small, above] at (2,5) {\textbf{Early Fusion}};
    \node[draw, fill=roxodna!30] (g1) at (0,4) {Genômica};
    \node[draw, fill=laranjacel!30] (t1) at (1.5,4) {Transcr.};
    \node[draw, fill=verdeacido!30] (p1) at (3,4) {Proteômica};
    \node[draw, fill=azulclaro!30] (m1) at (4.5,4) {Metabol.};
    \node[draw, fill=yellow!30, text width=1.5cm, align=center] (int1) at (2.25,2.5) {Integração\\Completa};
    \node[draw, fill=gray!30, text width=1.5cm, align=center] (mod1) at (2.25,1) {Modelo};

    \draw[->, thick] (g1) -- (int1);
    \draw[->, thick] (t1) -- (int1);
    \draw[->, thick] (p1) -- (int1);
    \draw[->, thick] (m1) -- (int1);
    \draw[->, thick] (int1) -- (mod1);

    % Intermediate fusion
    \node[font=\small, above] at (8,5) {\textbf{Intermediate Fusion}};
    \node[draw, fill=roxodna!30] (g2) at (6.5,4) {Genômica};
    \node[draw, fill=laranjacel!30] (t2) at (8,4) {Transcr.};
    \node[draw, fill=verdeacido!30] (p2) at (9.5,4) {Proteômica};
    \node[draw, fill=yellow!30] (int2a) at (7.25,2.8) {Módulo 1};
    \node[draw, fill=yellow!30] (int2b) at (9.25,2.8) {Módulo 2};
    \node[draw, fill=gray!30, text width=1.5cm, align=center] (mod2) at (8.25,1) {Modelo};

    \draw[->, thick] (g2) -- (int2a);
    \draw[->, thick] (t2) -- (int2a);
    \draw[->, thick] (p2) -- (int2b);
    \draw[->, thick] (int2a) -- (mod2);
    \draw[->, thick] (int2b) -- (mod2);

    % Late fusion
    \node[font=\small, above] at (13.5,5) {\textbf{Late Fusion}};
    \node[draw, fill=roxodna!30, text width=1.2cm, align=center] (g3) at (12,4) {Genômica};
    \node[draw, fill=laranjacel!30, text width=1.2cm, align=center] (t3) at (13.5,4) {Transcr.};
    \node[draw, fill=verdeacido!30, text width=1.2cm, align=center] (p3) at (15,4) {Proteômica};
    \node[draw, fill=yellow!30] (m3a) at (12,2.8) {Modelo 1};
    \node[draw, fill=yellow!30] (m3b) at (13.5,2.8) {Modelo 2};
    \node[draw, fill=yellow!30] (m3c) at (15,2.8) {Modelo 3};
    \node[draw, fill=gray!30, text width=1.5cm, align=center] (final) at (13.5,1) {Consenso};

    \draw[->, thick] (g3) -- (m3a);
    \draw[->, thick] (t3) -- (m3b);
    \draw[->, thick] (p3) -- (m3c);
    \draw[->, thick] (m3a) -- (final);
    \draw[->, thick] (m3b) -- (final);
    \draw[->, thick] (m3c) -- (final);
\end{tikzpicture}
\caption{Três estratégias de fusão de dados multi-ômicos. \textbf{Esquerda (Early Fusion):} todas as camadas são concatenadas em uma única matriz; o modelo recebe a matriz completa. \textbf{Centro (Intermediate Fusion):} cada camada é reduzida a representações latentes (módulos) --- e.g., MOFA, autoencoders ---, e os módulos são então combinados em um modelo único. \textbf{Direita (Late Fusion):} cada camada é modelada independentemente até o fim; o resultado final é um consenso entre os modelos (voto, média, regra bayesiana).}
\label{fig:11-fusao}
\end{figure}

\subsection{Métodos Baseados em Correlação}
\label{sub:image11-correlacao}

A abordagem mais simples é medir a correlação linear (Pearson) ou monotônica
(Spearman) entre features de camadas diferentes. Para duas séries $X$ e $Y$ em
condições pareadas:
\begin{equation}
\rho_{XY} = \frac{\operatorname{Cov}(X, Y)}{\sigma_X \sigma_Y}
\end{equation}

A correlação \emph{parcial} $\rho_{XY \mid Z}$ controla o efeito de um terceiro
conjunto de variáveis e é usada para distinguir correlação direta de
correlação mediada. A \emph{informação mútua} $I(X;Y)$ estende a análise a
relações não-monotônicas --- é mais robusta mas exige estimativa de densidades
em amostras pequenas. Já a \emph{correlação canônica} (CCA) busca combinações
lineares de duas variáveis multi-dimensionais que maximizem a correlação --- é
uma técnica clássica para detectar módulos compartilhados de variação entre
duas camadas.

Um alerta metodológico essencial: \emph{correlação não implica causalidade}. O
nível de mRNA e o nível de proteína correlacionam-se fortemente em escala
populacional, mas em muitos genes individuais essa correlação é fraca ou
inexistente (a abundância da proteína é regulada por degradação e tradução,
não apenas por síntese). As correlações encontradas em integração
multi-ômica são hipóteses a serem testadas por perturbações genéticas ou
farmacológicas, não afirmações finais sobre o sistema.

\subsection{Métodos Baseados em Redes}
\label{sub:image11-redes}

Cada camada pode ser representada como uma rede (\emph{grafo}) cujos nós são
features e cujas arestas representam interações. Redes \emph{multi-camada}
(\emph{multiplex}) usam o mesmo conjunto de nós em diferentes camadas --- mRNA,
proteína, metabólito --- com arestas intra-camada (interações dentro da mesma
camada ômica) e inter-camadas (interações entre features de camadas diferentes).
As análises mais frequentes são:

\begin{itemize}
\item \textbf{Centralidade} (grau, betweenness, closeness, eigenvector,
  PageRank): identifica nós centrais --- tipicamente reguladores ou
  hubs metabólicos.
\item \textbf{Modularidade} (Louvain, Infomap): identificação de comunidades
  densamente conectadas que provavelmente correspondem a unidades
  funcionais (complexos proteicos, módulos de co-expressão, vias metabólicas).
\item \textbf{Estrutura}: redes livres de escala, small-world, redes de
  Petersen, hierárquicas --- cada topologia com interpretação
  biológica específica.
\end{itemize}

\subsection{Aprendizado de Máquina Multi-ômico}
\label{sub:image11-ml}

A categoria mais ampla em integração é a de aprendizado de máquina. Métodos
supervisionados (Random Forest, SVM, redes neurais, Elastic Net para seleção
de features) usam dados rotulados --- tipicamente fenótipos clínicos --- para
treinar classificadores ou regressores sobre features multi-ômicas.
Métodos não-supervisionados (PCA, PLS, t-SNE, UMAP, NMF, autoencoders)
identificam estrutura sem rótulos --- ideal para exploração inicial ou
para descobrir subtipos populacionais não-antevistos. O método \emph{MOFA}
combina elementos de ambos --- uma fatoração supervisionada por uma
estrutura bayesiana esparsa, capaz de identificar fatores latentes compartilhados
e específicos de cada camada.

A escolha entre métodos não é, no entanto, um problema puramente técnico:
depende da estrutura dos dados, da presença ou ausência de rótulos
celulares, do tipo de integração desejada (correção de efeito de lote
versus modelagem probabilística), e do tamanho computacional admissível.
O estudo de referência para orientar essa escolha em contextos atlas-level
é o benchmarking conduzido por Luecken e colaboradores~\cite{luecken2022},
que comparou sistematicamente 68 combinações de método e pré-processamento
sobre 85 lotes de dados de scRNA-seq e scATAC-seq, totalizando mais de 1,2
milhão de células em 13 tarefas de integração atlas-level. Os autores
propuseram uma bateria padronizada de catorze métricas agrupadas em três
dimensões --- escalabilidade, capacidade de remoção de efeito de lote e
conservação de variação biológica --- e identificaram scANVI, Scanorama,
scVI e scGen como os métodos de melhor desempenho geral, com Harmony e
LIGER mostrando-se particularmente eficazes para integração de dados de
scATAC-seq. O pacote Python \texttt{scIB} e o pipeline de benchmarking
distribuídos pelos autores permitem que novos métodos ou novos conjuntos
de dados sejam avaliados sob o mesmo protocolo padronizado, oferecendo
um critério objetivo para a escolha de ferramentas em projetos concretos. A
Figura~\ref{fig:11-workflow} apresenta o pipeline típico de uma análise
multiômica: dos dados brutos heterogêneos (múltiplas plataformas) passa-se
por controle de qualidade, normalização, integração, análise estatística,
modelagem computacional e, finalmente, validação experimental --- com um
elo de retroalimentação (linha tracejada vermelha) que conecta a
validação de volta à modelagem.

\begin{figure}[htbp]
\centering
\begin{tikzpicture}[node distance=1.2cm, scale=0.75, transform shape,
    box/.style={rectangle, draw, text width=3cm, text centered, rounded corners, minimum height=0.8cm, font=\small}]

    \node[box, fill=roxodna!30] (raw) {Dados Brutos\\Multi-Ômicos};
    \node[box, fill=laranjacel!30, below of=raw] (qc) {Controle de\\Qualidade};
    \node[box, fill=verdeacido!30, below of=qc] (norm) {Normalização\\e Filtro};
    \node[box, fill=azulclaro!30, below of=norm] (int) {Integração\\de Dados};
    \node[box, fill=yellow!30, below of=int] (anal) {Análise\\Estatística};
    \node[box, fill=gray!30, below of=anal] (model) {Modelagem\\de Sistemas};
    \node[box, fill=green!30, below of=model] (val) {Validação\\Experimental};

    \draw[->, thick] (raw) -- (qc);
    \draw[->, thick] (qc) -- (norm);
    \draw[->, thick] (norm) -- (int);
    \draw[->, thick] (int) -- (anal);
    \draw[->, thick] (anal) -- (model);
    \draw[->, thick] (model) -- (val);

    % Feedback loop
    \draw[->, thick, dashed, red] (val) -| ++(2,0) |- (model);
\end{tikzpicture}
\caption{Pipeline típico de análise multiômica. A sequência canônica envolve sete etapas: (1) aquisição de dados brutos multi-ômicos (RNA-seq, proteômica, metabolômica, etc.); (2) controle de qualidade --- remoção de outliers, FastQC/MultiQC; (3) normalização e filtro (z-score, log, correção de efeito de lote); (4) integração de dados (correlação, redes, ML); (5) análise estatística (testes, FDR, enriquecimento); (6) modelagem de sistemas (ODEs, FBA, simulação); (7) validação experimental (qPCR, Western, ensaios enzimáticos). O elo de retroalimentação tracejado (vermelho) simboliza que resultados experimentais negativos podem levar a revisão do modelo.}
\label{fig:11-workflow}
\end{figure}

\begin{lstlisting}[language=Python, caption=Integracao multi-omica simples]
import pandas as pd
import numpy as np
from scipy import stats
import seaborn as sns
import matplotlib.pyplot as plt

trans = pd.read_csv('transcriptomics.csv', index_col=0)
prot  = pd.read_csv('proteomics.csv',     index_col=0)
metab = pd.read_csv('metabolomics.csv',   index_col=0)

def normalize_zscore(df):
    # Z-score por linha: cada gene centrado em zero
    return (df - df.mean()) / df.std()

t_n, p_n, m_n = map(normalize_zscore, [trans, prot, metab])

# Correlacao mRNA-proteina com filtro de p-valor
corr_mat = pd.DataFrame(index=t_n.index, columns=p_n.index)
for g in t_n.index:
    for pr in p_n.index:
        r, p = stats.pearsonr(t_n.loc[g], p_n.loc[pr])
        if p < 0.05: corr_mat.loc[g, pr] = r
        else:        corr_mat.loc[g, pr] = np.nan

plt.figure(figsize=(11, 9))
sns.heatmap(corr_mat.astype(float), cmap='RdBu_r', center=0,
            vmin=-1, vmax=1)
plt.title('Correlacao mRNA-Proteina (p < 0.05)')
plt.tight_layout(); plt.savefig('corrmatrix.png', dpi=150)
\end{lstlisting}

\subsection{Heterogeneidade dos Dados}
\label{sub:image11-heterogeneidade}

Cada camadaômica tem características próprias de distribuição e ruído. A
tabela a seguir resume os contrastes básicos:

\begin{itemize}
\item \textbf{Genômica}: aproximadamente 10$^4$ elementos, baixa faixa
  dinâmica (binário: heterozigoto/homozigoto/ausente), baixo ruído técnico.
\item \textbf{Transcriptômica}: 10$^4$--10$^7$ cópias por amostra, alta faixa
  dinâmica (fold-change de até 10$^4$), ruído moderado (variabilidade no
  número de reads).
\item \textbf{Proteômica}: 10$^4$ moléculas detectáveis, faixa dinâmica alta
  mas com \emph{missing values} frequentes, alto ruído técnico.
\item \textbf{Metabolômica}: 10$^2$--10$^3$ metabólitos identificados, faixa
  dinâmica muito alta (10$^9$), ruído alto.
\item \textbf{Fluxômica}: 10$^2$ fluxos estimados via $^{13}$C-MFA, faixa
  dinâmica média, ruído moderado.
\end{itemize}

A integração sem normalização apropriada é estatisticamente mal-condicionada
e gera conclusões espúrias. As práticas mais robustas incluem: z-score
intra-amostra ou inter-amostras, transformação Box-Cox para dados
contagem, correção de lote (\emph{batch correction}) por ComBat ou modelos
lineares mistos, e imputação de valores faltantes por k-NN ou métodos baseados
em expectativa-maximização --- sempre reportando a incerteza introduzida.


\section{Caso de Estudo: O Ciclo da Trealose em \textit{S. cerevisiae}}
\label{sec:11-trealose}

Para ilustrar com concretude o processo de integração multi-ômica, este
capítulo concentra-se em um sistema-modelo: o ciclo da trealose na levedura
\textit{Saccharomyces cerevisiae}. A escolha não é acidental --- a trealose é
um dos sistemas mais estudados da resposta ao estresse eucariótico, e os
mecanismos moleculares --- enzimáticos, regulatórios e alostéricos ---
estabelecidos para ele transferem-se com facilidade para sistemas mais
complexos, incluindo a proteção contra estresse em organismos multicelulares.

\subsection{Background Biológico: A Trealose como Protetor Celular}
\label{sub:image11-bio}

A trealose ($\alpha$-D-glicopiranosil-$\alpha$-D-glicopiranosídeo) é um
dissacarídeo não-redutor estabilizado por duas ligações glicosídicas
$\alpha{,}1 \to 1$, que protegem ambas as extremidades redutoras contra
reações de Maillard. Sua função primária na levedura é estabilizar membranas
e proteínas em estresses térmico, osmótico, oxidativo e de congelamento. Sob
essas condições, o dissacarídeo acumula até 20\% do peso seco celular; sob
condições favoráveis, é mobilizado como fonte de glicose.

Em aplicações industriais --- panificação, fermentação alcoólica,
biocombustíveis --- a trealose é valorizada por conferir robustez a
\textit{S. cerevisiae} em condições adversas. Em biotecnologia, modificar a
via é uma das estratégias mais antigas de engenharia de tolerância ao
estresse; em medicina, a trealose tem sido investigada como estabilizador de
proteínas terapêuticas e na preservação de órgãos para transplante.

\subsection{Via de Biossíntese e Degradação}
\label{sub:image11-via}

A síntese ocorre em duas etapas enzimáticas acopladas. \textit{TPS1}
(trealose-6-fosfato sintase) catalisa a condensação de UDP-glicose com
glicose-6-fosfato, formando trealose-6-fosfato (T6P) e UDP. \textit{TPS2}
(trealose-6-fosfato fosfatase) desfosforila T6P produzindo trealose livre.
A degradação é catalisada essencialmente por \textit{NTH1} (trealase
neutra), que hidrolisa trealose em duas glicoses livres.

A via tem três características regulatórias cruciais para a modelagem
integrada. \textbf{Primeira}: T6P é sintetizado a partir de G6P, e portanto
ocupa uma bifurcação entre a glicólise (G6P $\to$ frutose-6-fosfato) e a
biossíntese de trealose. O acúmulo de T6P sinaliza disponibilidade
energética suficiente para invest em reserva de estresse. \textbf{Segunda}:
o complexo TSC --- \textit{Tps1}, \textit{Tps2}, \textit{Tsl1}, \textit{Tps3}
--- tem cerca de 800 kDa e é altamente regulado por fosforilação. A
subunidade catalítica TPS1A é regulada por fosforilação por PKA, Snf1 e
Pho85. \textbf{Terceira}: a trealase neutra NTH1 é ativada por fosforilação
via PKA quando glicose retorna ao meio --- exatamente o momento em que a
trealose deve ser mobilizada. T6P também funciona como inibidor alostérico
de hexoquinase II e, em menor grau, regula a glicólise via PFK. A
Figura~\ref{fig:11-via-trealose} esquematiza a via completa: G6P e
UDP-Glc são convertidos em T6P pela enzima TPS1 (subunidade catalítica
do complexo TSC), T6P é desfosforilado por TPS2 produzindo trealose
livre, e a trealose é hidrolisada pela enzima NTH1 (trealase neutra) em
duas glicoses. A seta vermelha \emph{Estresse} $\to$ UDP-Glc indica a
ativação transcricional de TPS1 sob choque térmico; a seta azul
\emph{Recuperação} $\to$ Trealose indica a mobilização durante a
recuperação.

\begin{figure}[htbp]
\centering
\begin{tikzpicture}[scale=0.85, every node/.style={font=\small}]
    % Synthesis pathway
    \node[draw, circle, fill=azulclaro!30] (g6p) at (0,3) {G6P};
    \node[draw, circle, fill=azulclaro!30] (udpg) at (0,1.5) {UDP-Glc};
    \node[draw, circle, fill=azulclaro!30] (t6p) at (3,1.5) {T6P};
    \node[draw, circle, fill=laranjacel!50, text width=1.5cm, align=center] (tre) at (6,1.5) {Trealose};

    % Enzymes
    \node[draw, rectangle, fill=verdeacido!50, font=\tiny] (tps1) at (1.5,2.2) {TPS1};
    \node[draw, rectangle, fill=verdeacido!50, font=\tiny] (tps2) at (4.5,2.2) {TPS2};
    \node[draw, rectangle, fill=red!50, font=\tiny] (nth1) at (6,0) {NTH1};

    % Reactions
    \draw[->, ultra thick, azulescuro] (g6p) -- (udpg);
    \draw[->, ultra thick, azulescuro] (udpg) -- (t6p) node[above, midway, font=\tiny] {+ G6P};
    \draw[->, ultra thick, azulescuro] (t6p) -- (tre);
    \draw[->, ultra thick, red] (tre) -- (6,0.3) node[right, midway, font=\tiny] {2 Glc};

    % Enzyme labels
    \draw[->, dashed] (tps1) -- (1.5,1.5);
    \draw[->, dashed] (tps2) -- (4.5,1.5);
    \draw[->, dashed] (nth1) -- (6,0.75);

    % Regulation
    \node[draw, ellipse, fill=yellow!30, font=\tiny] (stress) at (-2,1.5) {Estresse};
    \draw[->, thick, red] (stress) -- (udpg) node[above, midway, font=\tiny] {$\uparrow$ síntese};

    \node[draw, ellipse, fill=green!30, font=\tiny] (normal) at (8,1.5) {Recuperação};
    \draw[->, thick, blue] (normal) -- (tre) node[above, midway, font=\tiny] {$\uparrow$ degradação};
\end{tikzpicture}
\caption{Via de biossíntese e degradação da trealose em \textit{Saccharomyces cerevisiae}. A síntese ocorre em duas etapas: TPS1 (trealose-6-fosfato sintase) catalisa UDP-Glc + G6P $\to$ T6P; TPS2 (T6P fosfatase) converte T6P em trealose livre. A degradação é catalisada por NTH1 (trealase neutra), que hidrolisa trealose em duas glicoses. A regulação é transcricional (estresse \emph{via} Msn2/4 ativa TPS1/TPS2), pós-traducional (PKA fosforila e ativa NTH1) e alostérica (T6P inibe PFK e hexoquinase II).}
\label{fig:11-via-trealose}
\end{figure}

\subsection{Dados Experimentais Multi-ômicos}
\label{sub:image11-dados}

O grupo de Nicolaas Pronk e colaboradores publicou, em 2014, um dos
conjuntos de dados multi-ômicos mais citados para o ciclo da trealose em
\textit{S. cerevisiae}. Sob choque térmico a 37°C, os autores mediram
simultaneamente RNA-seq (abundância de transcritos), proteômica
quantitativa (label-free) e metabolômica (LC-MS/MS) ao longo de 0, 15, 30,
60, 90 e 120 minutos. Os dados mostram uma resposta canônica: o
mRNA de \textit{TPS1} responde dentro de 15 minutos, atinge máximo em torno de
60 min e decai; a proteína TPS1 acumula com atraso de 30 min, alcança
plataforma em torno de 90 min; a trealose acumula de forma monotônica,
atingindo cerca de 130 mM no limite experimental.

Esse conjunto é ideal para discussão metodológica porque as três camadas
concorrem, mas com diferentes perfis temporais. A integração mostra que o
\emph{atraso transcrição--proteína} (~30 min) corresponde à meia-vida
canônica de tradução no citoplasma da levedura; o \emph{atraso
proteína--metabólito} adicional corresponde ao tempo de síntese e
reciclagem do substrato. Essas observações tornam-se \emph{predições} do
modelo quando colocadas em forma de EDOs e \emph{verificações} do modelo
quando confrontadas com a solução numérica.

\subsection{Construção de um Modelo Integrado}
\label{sub:image11-modelo}

A integração dos dados acima mencionados produz um modelo de quatro
variáveis em forma de EDOs:
\begin{align}
\frac{d[\mathrm{mRNA}_{TPS1}]}{dt} &= k_{\text{trans}}\, f(\text{estresse})
                                     - \delta_{\text{mRNA}}\,[\mathrm{mRNA}_{TPS1}] \\
\frac{d[TPS1]}{dt}               &= k_{\text{transl}}\,[\mathrm{mRNA}_{TPS1}]
                                     - \delta_{\text{prot}}\,[TPS1] \\
\frac{d[T6P]}{dt}                &= v_{TPS1} - v_{TPS2} \\
\frac{d[\mathrm{Trealose}]}{dt}  &= v_{TPS2} - v_{NTH1}
\end{align}

onde $f(\text{estresse})$ é uma função de Hill com coeficiente 4 e
constante de meio-ativação $K_d \approx 0{,}5$ (a intensidade de
estresse é a entrada do modelo). As velocidades seguem cinética Michaelis-Menten
para $v_{TPS2}$ e $v_{NTH1}$, e uma forma bi-substrato ordinária para
$v_{TPS1}$. As concentrações de G6P e UDP-glicose são tomadas constantes em
uma primeira aproximação --- o que está justificado pela observação
experimental de que seus níveis variam menos de 30\% no intervalo de
observação.

A estimação dos 11 parâmetros livres é feita por mínimos quadrados não
lineares com o método \texttt{scipy.optimize.least\_squares} (ou, com
incerteza populacional, MCMC via \texttt{PyMC3}). Restrições biológicas
($V_{\max}$, $K_m$, $\delta > 0$) e faixas operativas razoáveis para cada
parâmetro reduzem o espaço de busca e estabilizam o ajuste. A incorporação
explícita de dados multi-ômicos alivia a identificabilidade estrutural:
RNA-seq restringe $k_{\text{trans}}$ e $\delta_{\text{mRNA}}$; proteômica
restringe $k_{\text{transl}}$ e $\delta_{\text{prot}}$; metabolômica
restringe $V_{\max}^{TPS2}$, $V_{\max}^{NTH1}$ e seus $K_m$ associados.

A validação segue o ciclo de retro-alimentação mostrado no capítulo: o
modelo ajustado prevê, por exemplo, que superexpressão dez vezes maior de
TPS1 dobra a taxa de acúmulo de trealose, que deleção de NTH1 produz
trealose basal três vezes mais alta e que inibição de PKA aumenta a
produção basal em três vezes. Cada previsão é verificada
experimentalmente. O resultado é uma concordância tipicamente na faixa de
10\% entre previsão e observação --- evidência de que a integração
multi-ômica produz modelos com poder preditivo genuíno. A
Figura~\ref{fig:11-regulacao} detalha os três níveis regulatórios do
ciclo da trealose: regulação \textbf{transcricional} (estresse ativa o
fator Msn2/4 que induz a transcrição de TPS1/TPS2); regulação
\textbf{pós-traducional} (PKA fosforila e ativa NTH1 durante a
recuperação glicolítica); e regulação \textbf{alostérica} (T6P inibe
PFK, fechando o feedback entre a via de trealose e a glicólise).

\begin{figure}[htbp]
\centering
\begin{tikzpicture}[scale=0.8, every node/.style={font=\small}]
    % Transcriptional regulation
    \node[font=\footnotesize, above] at (2,4) {\textbf{Regulação Transcricional}};
    \node[draw, ellipse, fill=yellow!30] (stress) at (0,3) {Estresse};
    \node[draw, rectangle, fill=roxodna!30] (msn24) at (2,3) {Msn2/4};
    \node[draw, rectangle, fill=laranjacel!30] (tps1g) at (4,3) {\gene{TPS1}/\gene{TPS2}};

    \draw[->, thick] (stress) -- (msn24);
    \draw[->, thick] (msn24) -- (tps1g);

    % Post-translational regulation
    \node[font=\footnotesize, above] at (2,1.5) {\textbf{Regulação Pós-Traducional}};
    \node[draw, ellipse, fill=verdeacido!30] (glucose) at (0,0.5) {Glicose};
    \node[draw, rectangle, fill=azulclaro!30] (camp) at (2,0.5) {cAMP/PKA};
    \node[draw, rectangle, fill=yellow!30] (nth1p) at (4,0.5) {NTH1-P};

    \draw[->, thick] (glucose) -- (camp);
    \draw[->, thick] (camp) -- (nth1p) node[above, midway, font=\tiny] {fosforilação};

    % Allosteric regulation
    \node[font=\footnotesize, above] at (9,4) {\textbf{Regulação Alostérica}};
    \node[draw, circle, fill=azulclaro!30] (g6p2) at (7.5,3) {G6P};
    \node[draw, circle, fill=laranjacel!50] (t6p2) at (9,3) {T6P};
    \node[draw, rectangle, fill=verdeacido!30] (pfk) at (10.5,3) {PFK};

    \draw[->, thick, red] (t6p2) -- (pfk) node[above, midway, font=\tiny] {inibição};
    \draw[->, thick, blue] (g6p2) -- (pfk) node[above, midway, font=\tiny] {ativação};
\end{tikzpicture}
\caption{Os três níveis de regulação do ciclo da trealose. \textbf{Transcricional:} estresse térmico ativa o fator de transcrição Msn2/4, que se liga ao \emph{stress response element} (STRE) nos promotores de \gene{TPS1} e \gene{TPS2}, induzindo sua transcrição. \textbf{Pós-traducional:} retorno da glicose ativa adenilato ciclase $\to$ cAMP $\to$ PKA, que fosforila NTH1 e a ativa, mobilizando trealose. \textbf{Alostérica:} T6P acumulado inibe PFK (glicólise), fechando um feedback negativo --- alta trealose desacelera a glicólise, preservando G6P.}
\label{fig:11-regulacao}
\end{figure}

\begin{lstlisting}[language=Python, caption=Modelo integrado do ciclo da trealose]
import numpy as np
from scipy.integrate import odeint
import matplotlib.pyplot as plt

def trehalose_model(y, t, params):
    mRNA, TPS1, T6P, Tre = y
    (k_trans, delta_mRNA, k_transl, delta_prot,
     V_TPS1, Km_G6P, V_TPS2, Km_T6P,
     V_NTH1, Km_Tre, stress) = params

    # Substratos aproximados como constantes
    G6P, UDP_Glc = 5.0, 2.0   # mM

    # Regulacao transcricional - Hill
    f_stress = stress**4 / (0.5**4 + stress**4)

    v_TPS1 = (V_TPS1 * TPS1 * G6P * UDP_Glc) / ((Km_G6P + G6P) * (1 + UDP_Glc))
    v_TPS2 = (V_TPS2 * T6P) / (Km_T6P + T6P)
    v_NTH1 = (V_NTH1 * Tre) / (Km_Tre + Tre)

    dmRNA  = k_trans * f_stress - delta_mRNA * mRNA
    dTPS1  = k_transl * mRNA - delta_prot * TPS1
    dT6P   = v_TPS1 - v_TPS2
    dTre   = v_TPS2 - v_NTH1

    return [dmRNA, dTPS1, dT6P, dTre]

params = [5.0, 0.10, 2.0, 0.05,    # transcricao e traducao
          10.0, 0.5, 15.0, 0.2,    # TPS1, TPS2
          8.0, 10.0, 1.0]          # NTH1, stress

t = np.linspace(0, 120, 200)
y0 = [1.0, 1.0, 0.1, 5.0]
sol = odeint(trehalose_model, y0, t, args=(params,))

fig, axes = plt.subplots(2, 2, figsize=(12, 9))
labels = ['mRNA TPS1', 'Proteina TPS1', 'T6P', 'Trealose (mM)']
for i, (ax, lab) in enumerate(zip(axes.flat, labels)):
    ax.plot(t, sol[:, i], 'b-', lw=2, label=labels[i])
    ax.set_xlabel('Tempo (min)'); ax.legend(); ax.grid(alpha=0.3)

# Dados experimentais: trealose observada
t_exp = np.array([0, 15, 30, 60, 90, 120])
tre_exp = np.array([5, 8, 25, 85, 120, 130])
axes[1, 1].plot(t_exp, tre_exp, 'ko', ms=8, label='observado')
axes[1, 1].legend()
plt.tight_layout(); plt.savefig('trealose_fit.png', dpi=150)
\end{lstlisting}


\section{Análise de Vias e Redes Integradas}
\label{sec:11-vias-redes}

Integrar dados multi-ômicos tem um objetivo final que vai além de gerar
\emph{listas de genes}: é produzir \emph{interpretações biológicas} das
alterações. Duas famílias de ferramentas dominam essa interpretação --- a
análise de enriquecimento (estática) e a reconstrução de redes (dinâmica)
--- e ambas evoluíram significativamente nos últimos anos para acomodar
dados de múltiplas camadas.

\subsection{Análise de Enriquecimento de Vias (ORA)}
\label{sub:image11-ora}

A análise de enriquecimento, em sua forma mais simples (Over-Representation
Analysis, ORA), responde à pergunta: o conjunto de genes diferencialmente
expressos em minha condição está \emph{mais presente} do que o acaso
esperaria em uma via biológica específica? Formalmente, o teste é o
hipergeométrico:
\begin{equation}
p = \frac{\binom{K}{k}\,\binom{N - K}{n - k}}{\binom{N}{n}}
\end{equation}
onde $N$ é o número total de genes no genoma, $K$ o número de genes na via
de interesse, $n$ o número de genes no conjunto em teste, e $k$ o número
de genes na interseção. O método é direto e dezenas de ferramentas o
implementam --- DAVID, Enrichr, g:Profiler, clusterProfiler (em R). Sua
desvantagem fundamental é a dependência de um \emph{threshold} arbitrário
para declarar genes como diferencialmente expressos.

\subsection{Gene Set Enrichment Analysis (GSEA)}
\label{sub:image11-gsea}

O GSEA, introduzido por Subramanian et al.\ (2005), substitui o threshold
arbitrário por um \emph{ranking} contínuo de genes: em vez de perguntar
"este gene está diferencialmente expresso?", pergunta "o conjunto de genes
desta via está deslocado para cima ou para baixo no ranking de fold-change?".

O procedimento é simples em descrição. Calcular um ranking ordenado
$\{r_1, r_2, \dots, r_N\}$ dos genes --- usualmente pela estatística-t ou
pelo log-fold-change --- e calcular o \emph{Enrichment Score} (ES) sobre
cada via: à medida que se percorre o ranking, soma-se $\sqrt{\tfrac{N - S}{S}}$
quando o gene pertence à via e subtrai-se $\sqrt{\tfrac{S}{N - S}}$
quando não pertence. O ES é o máximo da soma parcial; é positivo se a via
tende a estar no topo do ranking, negativo se tende a estar na base.

A significância do ES é avaliada por \emph{permutação dos fenótipos}: as
amostras (não os genes) são embaralhadas várias centenas de vezes, o
procedimento é repetido, e a posição do ES observado na distribuição nula
fornece um p-valor empírico. A correção para múltiplos testes usa FDR
(\emph{false discovery rate}) de Benjamini-Hochberg.

A vantagem do GSEA em relação ao ORA é a capacidade de detectar mudanças
\emph{coordenadas mas sutis}, mesmo quando genes individuais não cruzam o
threshold de significância. Para estudos com poucas amostras, no entanto, a
permutação é instável e o método perde poder --- neste caso, ferramentas
como \texttt{fgsea} em R, que aproximam o p-valor por distribuição teórica,
são mais robustas. A Figura~\ref{fig:11-gsea} apresenta um exemplo
típico do gráfico de enriquecimento: o eixo horizontal lista os genes
ordenados pelo ranking de fold-change (de mais superexpressos à esquerda
para mais subexpressos à direita); a curva azul mostra o \emph{enrichment
score} (ES) acumulado à medida que genes da via são encontrados; o
pico (ES$_{\max}$, ponto vermelho) indica o ponto de máximo
enriquecimento.

\begin{figure}[htbp]
\centering
\begin{tikzpicture}[scale=0.6]
    % GSEA plot
    \draw[->] (0,0) -- (6,0) node[right] {Genes (ranked)};
    \draw[->] (0,0) -- (0,3) node[above] {ES};

    % Enrichment curve
    \draw[thick, azulescuro, domain=0:6, samples=50] plot (\x, {1.5 + 0.8*sin(180*\x/3)});

    % Gene set hits
    \foreach \x in {0.5, 1.0, 1.2, 1.5, 1.8, 2.2} {
        \draw[red, thick] (\x,0) -- (\x,0.3);
    }

    % Labels
    \node[below, font=\tiny] at (1,0) {Alta expr.};
    \node[below, font=\tiny] at (5,0) {Baixa expr.};
    \node[left, font=\tiny] at (0,2.3) {ES$_{max}$};

    % Enrichment peak
    \draw[dashed] (0,2.3) -- (1.8,2.3);
    \draw[dashed] (1.8,0) -- (1.8,2.3);
    \node[circle, fill=red, inner sep=1pt] at (1.8,2.3) {};
\end{tikzpicture}
\caption{Gráfico típico de Gene Set Enrichment Analysis (GSEA). O eixo horizontal representa os genes ordenados pelo ranking de fold-change (de alta expressão à esquerda para baixa expressão à direita); o eixo vertical é o \emph{enrichment score} (ES) acumulado. A curva azulescura é o ES à medida que genes da via sendo testada são encontrados; os traços vermelhos no eixo horizontal marcam as posições dos \emph{hits}. O pico ES$_{\max}$ (ponto vermelho, ~$1{,}8$) indica o ponto de máximo enriquecimento --- ES positivo significa via ativada; ES negativo significa via reprimida.}
\label{fig:11-gsea}
\end{figure}

\subsection{Reconstrução de Redes a partir de Multi-ômicos}
\label{sub:image11-reconstrucao}

A construção de redes a partir de dados multi-ômicos usa múltiplas
estratégias, cada uma com suas hipóteses embutidas. A correlação (Pearson)
identifica relações lineares; a correlação parcial controla co-variáveis
confundidoras pela fórmula:
\begin{equation}
\rho_{XY \mid Z} = \frac{\rho_{XY} - \rho_{XZ}\rho_{YZ}}
                        {\sqrt{(1 - \rho_{XZ}^2)(1 - \rho_{YZ}^2)}}
\end{equation}

Informação mútua detecta relações não-monotônicas:
\begin{equation}
I(X;Y) = \sum_{x, y} p(x, y)\, \log \frac{p(x, y)}{p(x)\, p(y)}
\end{equation}

Graphical LASSO (gLASSO) infere grafos gaussianos esparsos a partir de
correlações amostrais, penalizando o número de arestas. Redes Bayesianas
dinâmicas aprendem grafos causais em séries temporais, com a vantagem
de distinguir causa de efeito pela ordem temporal --- são computacionalmente
caras mas produzem as reconstruções mais fidedignas.

\subsection{Detecção de Módulos e Hubs}
\label{sub:image11-modulos}

Uma vez obtida a rede, a análise subsequente busca módulos funcionais ---
grupos de nós densamente conectados que provavelmente representam unidades
biológicas. O algoritmo \emph{Louvain} maximiza a modularidade $Q$:
\begin{equation}
Q = \frac{1}{2m}\sum_{i, j} \left[ A_{ij} - \frac{k_i k_j}{2m} \right]
                           \delta(c_i, c_j)
\end{equation}
onde $A_{ij}$ é a matriz de adjacência, $k_i$ é o grau do nó $i$ e $c_i$ é
a comunidade atribuída. Valores típicos de $Q$ acima de $0{,}3$ indicam
estrutura modular significativa.

A centralidade é complementar e identifica \emph{hubs} --- nós centrais que
provavelmente representam reguladores. Centralidade de grau mede o número
direto de conexões; betweenness mede se o nó é um gargalo para fluxos na
rede; closeness mede proximidade média; eigenvector favorece vizinhos que
já são importantes (e ignora nós periféricos conectados a centrais).

A integração multi-ômica adiciona uma camada adicional: módulos que
co-variam consistentemente entre múltiplas camadas ômicas são fortes
candidatos a \emph{unidades regulatórias reais}, e seus hubs são candidatos
a reguladores mestres da resposta em estudo. Em leveduras, a abordagem
identificou o fator de transcrição Msn2/4 como hub central da resposta a
choque térmico --- confirmação que valida toda a estratégia de análise. A
Figura~\ref{fig:11-modulos} ilustra o resultado típico de uma análise
de modularidade (algoritmo de Louvain) aplicada a uma rede de
co-expressão: três comunidades densamente conectadas (módulos), cada
uma com sua cor característica, separadas por arestas inter-módulo
(linhas tracejadas) esparsas --- a assinatura de uma rede com estrutura
modular.

\begin{figure}[htbp]
\centering
\begin{tikzpicture}[scale=0.6]
    % Network with modules
    % Module 1
    \node[draw, circle, fill=azulclaro!50, font=\tiny] (a1) at (0,2) {G1};
    \node[draw, circle, fill=azulclaro!50, font=\tiny] (a2) at (1,2.5) {G2};
    \node[draw, circle, fill=azulclaro!50, font=\tiny] (a3) at (1,1.5) {G3};

    \draw[thick] (a1) -- (a2);
    \draw[thick] (a1) -- (a3);
    \draw[thick] (a2) -- (a3);

    % Module 2
    \node[draw, circle, fill=verdeacido!50, font=\tiny] (b1) at (3,2) {G4};
    \node[draw, circle, fill=verdeacido!50, font=\tiny] (b2) at (4,2.5) {G5};
    \node[draw, circle, fill=verdeacido!50, font=\tiny] (b3) at (4,1.5) {G6};

    \draw[thick] (b1) -- (b2);
    \draw[thick] (b1) -- (b3);
    \draw[thick] (b2) -- (b3);

    % Module 3
    \node[draw, circle, fill=laranjacel!50, font=\tiny] (c1) at (2,0) {G7};
    \node[draw, circle, fill=laranjacel!50, font=\tiny] (c2) at (3,0) {G8};

    \draw[thick] (c1) -- (c2);

    % Inter-module edges
    \draw[dashed, thin] (a3) -- (b1);
    \draw[dashed, thin] (b3) -- (c1);

    % Labels
    \node[below, font=\tiny] at (0.5,3) {Módulo 1};
    \node[below, font=\tiny] at (3.5,3) {Módulo 2};
    \node[below, font=\tiny] at (2.5,-0.5) {Módulo 3};
\end{tikzpicture}
\caption{Detecção de módulos em rede de co-expressão. Oito genes (G1--G8) se organizam em três comunidades densamente conectadas (Módulo 1 --- azul-claro, Módulo 2 --- verde, Módulo 3 --- laranja), detectadas por maximização da modularidade $Q$ (algoritmo de Louvain). Arestas sólidas representam conexões intra-módulo; arestas tracejadas representam conexões inter-módulo esparsas --- assinatura de estrutura modular significativa ($Q > 0{,}3$).}
\label{fig:11-modulos}
\end{figure}

\begin{lstlisting}[language=Python, caption=Rede de co-expressao e deteccao de hubs]
import networkx as nx
import pandas as pd
import matplotlib.pyplot as plt
from scipy.stats import pearsonr

def build_correlation_network(data, threshold=0.7, alpha=0.05):
    # Construir grafo com threshold em |correlacao| e p-valor
    G = nx.Graph()
    for g in data.index: G.add_node(g)
    for i, g1 in enumerate(data.index):
        for g2 in data.index[i+1:]:
            r, p = pearsonr(data.loc[g1], data.loc[g2])
            if abs(r) > threshold and p < alpha:
                G.add_edge(g1, g2, weight=abs(r))
    return G

def analyze_network(G):
    deg_c = nx.degree_centrality(G)
    btw_c = nx.betweenness_centrality(G)
    clo_c = nx.closeness_centrality(G)
    threshold = sorted(deg_c.values(), reverse=True)[max(1, len(G)//10)]
    hubs = [n for n, c in deg_c.items() if c >= threshold]
    return deg_c, btw_c, clo_c, hubs

# Exemplo: dados de co-expressao
data = pd.read_csv('expression_ts.csv', index_col=0)
G = build_correlation_network(data, threshold=0.8)
deg, btw, clo, hubs = analyze_network(G)

print(f"Rede: {G.number_of_nodes()} nos, {G.number_of_edges()} arestas")
print(f"Top hubs: {hubs[:10]}")

# Plotar
pos = nx.spring_layout(G, k=0.5, seed=42)
node_color = ['red' if n in hubs else 'lightblue' for n in G.nodes()]
nx.draw(G, pos, node_color=node_color, node_size=80,
        edge_color='gray', with_labels=False, alpha=0.7)
plt.title('Rede de co-expressao (hubs em vermelho)')
plt.savefig('network.png', dpi=150)
\end{lstlisting}

A Figura~\ref{fig:11-redes-din} ilustra a evolução temporal de uma
rede sob estresse térmico: em $t=0$ a rede é esparsa (poucos nós, arestas
finas); em $t = 30$ min aparecem novas conexões; em $t = 60$ min a
rede se reorganiza completamente --- alguns nós aumentam de tamanho
(maior \emph{degree}, indicando que se tornaram hubs regulatórios).
A análise dinâmica de redes capta essa sequência de eventos.

\begin{figure}[htbp]
\centering
\begin{tikzpicture}[scale=0.6]
    % Time series of networks
    \node[font=\tiny] at (0,3) {t = 0};
    \node[draw, circle, fill=azulclaro!50, minimum size=0.3cm] (t0a) at (0,2) {};
    \node[draw, circle, fill=azulclaro!50, minimum size=0.3cm] (t0b) at (0.8,2) {};
    \node[draw, circle, fill=azulclaro!50, minimum size=0.3cm] (t0c) at (0.4,1.2) {};
    \draw[thick] (t0a) -- (t0b);

    \node[font=\tiny] at (2.5,3) {t = 30 min};
    \node[draw, circle, fill=laranjacel!50, minimum size=0.4cm] (t1a) at (2.5,2) {};
    \node[draw, circle, fill=laranjacel!50, minimum size=0.3cm] (t1b) at (3.3,2) {};
    \node[draw, circle, fill=laranjacel!50, minimum size=0.3cm] (t1c) at (2.9,1.2) {};
    \draw[thick] (t1a) -- (t1b);
    \draw[thick] (t1a) -- (t1c);
    \draw[thick] (t1b) -- (t1c);

    \node[font=\tiny] at (5,3) {t = 60 min};
    \node[draw, circle, fill=verdeacido!50, minimum size=0.6cm] (t2a) at (5,2) {};
    \node[draw, circle, fill=verdeacido!50, minimum size=0.5cm] (t2b) at (5.8,2) {};
    \node[draw, circle, fill=verdeacido!50, minimum size=0.4cm] (t2c) at (5.4,1.2) {};
    \draw[ultra thick] (t2a) -- (t2b);
    \draw[ultra thick] (t2a) -- (t2c);
    \draw[thick] (t2b) -- (t2c);

    % Arrows showing time
    \draw[->, thick] (1,1.5) -- (1.8,1.5);
    \draw[->, thick] (3.5,1.5) -- (4.3,1.5);

    \node[below, font=\tiny, text width=2cm, align=center] at (3,-0.2) {Resposta ao estresse térmico};
\end{tikzpicture}
\caption{Evolução temporal de uma rede gênica sob estresse térmico. Três instantâneos ($t = 0$, $t = 30$ min, $t = 60$ min) mostram a reorganização da topologia: o número de arestas cresce com o tempo (resposta transcricional), e alguns nós aumentam de tamanho --- tornando-se \emph{hubs} --- à medida que o sistema se adapta. A análise dinâmica de redes (sliding window, dynamic Bayesian networks, graphical Granger causality) permite identificar a sequência temporal de eventos regulatórios em resposta à perturbação.}
\label{fig:11-redes-din}
\end{figure}

\section{Modelagem de Sistemas Integrados}
\label{sec:11-modelagem}

A integração multi-ômica não termina na estatística: ela alimenta
\emph{modelos de sistemas}. Nesta seção discutimos os três formalismos
dominantes --- modelos em escala genômica (GEMs) baseados em restrições,
modelos cinéticos baseados em EDOs, e modelos híbridos que combinam os
anteriores ---, dando ênfase especial ao uso de dados multi-ômicos como
\emph{restrições} ou \emph{entradas parametrizadas} dos modelos.

\subsection{Modelos em Escala Genômica com Restrições Ômicas}
\label{sub:image11-gem}

\begin{definicaoenv}[Modelo em escala genômica (GEM)]
Reconstrução matemática do metabolismo completo de um organismo, incluindo
todas as reações conhecidas com estequiometria, modificadores enzimáticos e
associações com genes. Para \textit{S. cerevisiae}, o modelo mais usado é
o iMM904 (~1500 genes, ~2000 reações). Para humanos, modelos como o Recon3D
oferecem cobertura metabolicamente completa (~13 000 reações).
\end{definicaoenv}

A integração de dados transcriptômicos em GEMs segue três linhas
principais. \textbf{E-Flux} afirma que a capacidade de fluxo de cada reação
é proporcional ao nível de expressão de seus genes associados: $v_{\max,i}
\propto \text{expressão}_i$. O modelo original é resolvido por FBA padrão,
mas com restrições modificadas. \textbf{GIMME} minimiza o desacordo entre
fluxos preditos e dados de expressão binarizados: para reações com
evidência de alta expressão, $v_i$ deve aproximar-se de um valor-alvo
$v_i^{\text{target}}$; para reações com baixa evidência, o custo é zero.
\textbf{iMAT} usa dados categóricos (on/off) e maximiza o número de
reações consistentes com a observação.

A análise central em todos esses métodos é a Análise por Balanceamento de
Fluxos (FBA --- Flux Balance Analysis), que resolve o problema de
programação linear:
\begin{align}
\max_v \quad & c^{\top} v \\
\text{sujeito a} \quad & S v = 0 \\
                       & \mathrm{lb} \leq v \leq \mathrm{ub}
\end{align}
onde $S$ é a matriz estequiométrica, $v$ o vetor de fluxos, $c$ a função
objetivo (por exemplo, maximização de biomassa) e os limites lb, ub são as
restrições. Quando esses limites são ajustados por dados de expressão, a
solução ótima representa o fluxo metabolicamente consistente com a
condição celular medida.

O COBRApy é a implementação padrão em Python. O trecho de código abaixo
ilustra a aplicação de restrições de expressão sobre o modelo iMM904 da
levedura:

\begin{lstlisting}[language=Python, caption=FBA integrado com transcriptomica]
import cobra
import pandas as pd

model = cobra.io.read_sbml_model('iMM904.xml')
expr  = pd.read_csv('expression_stress.csv', index_col=0)
expr_norm = expr / expr.max()

def integrate_expression(model, expression, pct_block=25):
    # Bloqueia reacoes em baixa expressao; escala as demais
    for r in model.reactions:
        genes = list(r.genes)
        if not genes: continue
        values = [float(expression.loc[g.id].mean())
                  for g in genes if g.id in expression.index]
        if not values: continue
        avg = sum(values) / len(values)
        if avg * 100 < pct_block:
            r.upper_bound = 0
            r.lower_bound = 0
        else:
            ub0 = r.upper_bound
            if ub0 > 0:
                r.upper_bound = max(ub0 * avg, 1e-6)

integrate_expression(model, expr_norm)
sol = model.optimize()
print(f"Taxa de crescimento com restricao: "
      f"{sol.objective_value:.3f}/h")
\end{lstlisting}

A Figura~\ref{fig:11-fba-restricao} ilustra conceitualmente o efeito de
restrições de expressão em FBA. O \emph{feasible space} da FBA padrão
(região azul-clara grande) contém todas as soluções de fluxo
viáveis. As restrições de expressão --- tipicamente limitando reações
associadas a genes pouco expressos e escalando as demais pela média de
expressão --- reduzem esse espaço (região verde sobreposta), movendo o
ponto ótimo para uma localização metabolicamente mais realista.

\begin{figure}[htbp]
\centering
\begin{tikzpicture}[scale=0.6]
    % FBA concept
    \node[font=\tiny, above] at (2,4) {\textbf{Espaço de Soluções}};

    % Feasible region
    \fill[azulclaro!20] (0,0) -- (4,0) -- (4,3) -- (0,3) -- cycle;

    % Expression constraints
    \fill[verdeacido!40] (0.5,0.5) -- (3.5,0.5) -- (3.5,2.5) -- (0.5,2.5) -- cycle;

    % Optimal solution
    \node[circle, fill=red, inner sep=2pt, label=right:{Ótimo}] at (3,2) {};

    % Axes
    \draw[->] (0,0) -- (4.5,0) node[right, font=\tiny] {$v_1$};
    \draw[->] (0,0) -- (0,3.5) node[above, font=\tiny] {$v_2$};

    % Labels
    \node[below, font=\tiny, text width=2cm, align=center] at (2,-0.5) {FBA padrão};
    \node[below, font=\tiny, text width=2cm, align=center] at (2,1.5) {Com expressão};
\end{tikzpicture}
\caption{Esquema conceitual de FBA com restrições de expressão. O espaço de soluções viáveis do problema de FBA (região azul-clara grande) é definido por todas as restrições estequiométricas e termodinâmicas. A introdução de restrições baseadas em dados transcriptômicos (E-Flux, GIMME, iMAT) reduz esse espaço (região verde) --- bloqueando reações associadas a genes pouco expressos e escalando as demais pela média de expressão. O ponto ótimo (vermelho) se move para uma solução metabolicamente mais realista, condizente com a condição celular medida.}
\label{fig:11-fba-restricao}
\end{figure}

\subsection{Modelos Cinéticos com Proteômica e Metabolômica}
\label{sub:image11-cineticos}

Quando se deseja dinâmica --- e não apenas estado estacionário --- a
abordagem natural são modelos de EDOs. A integração multi-ômica entra de
duas formas. Primeiro, dados proteômicos quantitativos definem
concentrações absolutas das enzimas $[E_i]$, ao invés de tratá-las como
parâmetros livres. Segundo, dados metabolômicos quantitativos
substituem valores de $[S_i]$, $[P_i]$ que de outro modo precisariam ser
medidos em tempo real.

Para cada reação $i$ no modelo cinético, a expressão genérica é:
\begin{equation}
v_i = [E_i] \, k_{\mathrm{cat},i} \, \frac{[S_i]}{K_m^S + [S_i]}
                     \left( 1 - \frac{[P_i]}{[S_i] K_{\mathrm{eq},i}} \right)
\end{equation}
onde o fator final $(1 - [P_i]/([S_i] K_{\mathrm{eq},i}))$ é a correção
termodinâmica. O sistema de EDOs é então:
\begin{equation}
\frac{d[S_i]}{dt} = \sum_j s_{ij} v_j
\end{equation}
onde $s_{ij}$ é o coeficiente estequiométrico.

A combinação de quatro fontes de dados (RNA-seq, proteômica, metabolômica
e dados termodinâmicos $\Delta G^0$) produz um modelo em que os parâmetros
livre restantes --- principalmente $k_{\mathrm{cat},i}$ --- podem ser
obtidos por ajuste ou por catálogos como \texttt{SabioRK} e \texttt{BRENDA}.

\subsection{Modelagem Híbrida}
\label{sub:image11-hibridos}

\begin{definicaoenv}[Modelo híbrido]
Modelo que combina diferentes formalismos matemáticos em camadas distintas
do sistema --- por exemplo, FBA para metabolismo central + EDOs para
regulação transcricional + Gillespie para estocasticidade de baixa
abundância. Modelos híbridos capturam múltiplas escalas temporais e
distribuem a complexidade computacional de forma realista.
\end{definicaoenv}

A escolha do formalismo por camada segue a escala temporal: segundos a
minutos --- EDOs (metabolismo e sinalização); horas --- EDOs
(transcrição, tradução); steady-state --- FBA (metabolismo para escala
genômica); raros eventos estocásticos --- Gillespie (número pequeno de
moléculas). Em \textit{S. cerevisiae}, modelos híbridos foram construídos
para integrar metabolismo, regulação gênica e controle de divisão celular
--- e demonstraram capacidade preditiva para

fenótipos de mutantes que escapariam a modelos mono-formalismos.

A integração de dados multi-ômicos ocorre em cada camada:
dados transcriptômicos alimentam o componente regulatório, dados
proteômicos alimentam o componente enzimático, metabolômicos alimentam o
componente metabólico de baixa escala. A grande desvantagem é a
identificabilidade: quanto mais camadas, mais difícil garantir que existe
uma solução única para os parâmetros. Estratégias de redução (dados
multi-ômicos fixando variáveis, hierarquias bayesianas, regressão
esparsa) são essenciais para manter a parcimônia.

A validação experimental completa o ciclo. As previsões do modelo híbrido
--- taxa de acúmulo de trealose, resposta transcricional após pulso de
glicose, viabilidade de mutantes sob estresse combinado --- devem ser
confrontadas com dados independentes antes de considerar o modelo útil.
Estudos recentes em biologia de sistemas de leveduras usam essa estratégia
com sucesso --- por exemplo, o modelo de \textit{S. cerevisiae} de
Dinh et al.\ (2022) integra mais de 50 condições de estresse e dezenas de
milhares de medidas em um único modelo coerente. A
Figura~\ref{fig:11-hibrido} ilustra a estrutura de um modelo híbrido
multi-escala: três camadas (regulação gênica via EDOs, tradução
estocástica via Gillespie, metabolismo central via FBA em estado
estacionário) acopladas em série, cada uma com sua escala temporal
característica (min-h, seg-min, steady-state, respectivamente).

\begin{figure}[htbp]
\centering
\begin{tikzpicture}[scale=0.7]
    % Hybrid model structure
    \node[draw, rectangle, fill=roxodna!30, text width=2cm, align=center, font=\tiny] at (0,3) {Regulação\\Gênica\\(ODE)};
    \node[draw, rectangle, fill=laranjacel!30, text width=2cm, align=center, font=\tiny] at (0,1.5) {Tradução\\(Estocástico)};
    \node[draw, rectangle, fill=verdeacido!30, text width=2cm, align=center, font=\tiny] at (0,0) {Metabolismo\\(FBA)};

    % Connections
    \draw[<->, thick] (0,2.7) -- (0,1.8);
    \draw[<->, thick] (0,1.2) -- (0,0.3);

    % Time scales
    \node[right, font=\tiny] at (2.5,3) {Escala: min-h};
    \node[right, font=\tiny] at (2.5,1.5) {Escala: seg-min};
    \node[right, font=\tiny] at (2.5,0) {Escala: steady-state};
\end{tikzpicture}
\caption{Esquema de modelo híbrido multi-escala. Três camadas de modelagem com diferentes formalismos --- EDOs para regulação gênica (escala horas), Gillespie para tradução estocástica (escala segundos-minutos), FBA para metabolismo central (estado estacionário) --- acopladas em série. Cada camada usa o formalismo mais adequado à sua escala temporal; os parâmetros das camadas são alimentados por dados multi-ômicos (RNA-seq $\to$ EDOs, scRNA-seq $\to$ Gillespie, metabolômica $\to$ FBA). O resultado é um modelo com poder preditivo em múltiplas escalas, que pode ser validado contra dados independentes em cada nível.}
\label{fig:11-hibrido}
\end{figure}

\section{Exercícios}
\label{sec:11-exercicios}

\begin{exercicio}
\textbf{Correlação multiômica.} Em um experimento de choque térmico em
levedura, foram medidos os níveis de mRNA e proteína para cinco genes:

\begin{center}
\begin{tabular}{lcc}
\toprule
\textbf{Gene} & \textbf{log$_2$ mRNA} & \textbf{log$_2$ proteína} \\
\midrule
\textit{TPS1} & 3{,}2 & 1{,}8 \\
\textit{HSP104} & 4{,}5 & 2{,}1 \\
\textit{CTT1} & 2{,}8 & 1{,}5 \\
\textit{SSA1} & 3{,}7 & 2{,}0 \\
\textit{GLK1} & 0{,}5 & 0{,}3 \\
\bottomrule
\end{tabular}
\end{center}

(a) Calcule a correlação de Pearson entre as duas séries. (b) Avalie sua
significância estatística. (c) Os genes estão em uma faixa de variação
suficiente para validar o resultado? (d) Interprete biologicamente, em
particular o que pode estar acontecendo com \textit{GLK1}.
\end{exercicio}

\begin{exercicio}
\textbf{Teste hipergeométrico.} Em um experimento de estresse oxidativo em
levedura, 150 genes foram diferencialmente expressos em um genoma de 6000
genes; a via de Resposta ao Estresse Oxidativo contém 80 genes, dos quais
25 estão na lista diferencial. (a) Calcule o p-valor pelo teste
hipergeométrico. (b) A via está significativamente enriquecida ao nível
$\alpha = 0{,}05$? (c) Aplique correção de Bonferroni para 300 vias testadas
e discuta a mudança de conclusão. (d) Por que FDR (Benjamini-Hochberg) é
uma alternativa preferível?
\end{exercicio}

\begin{exercicio}
\textbf{GSEA vs.\ ORA.} (a) Explique conceitualmente por que o GSEA detecta
mudanças coordenadas sutis que o ORA perde. (b) Em quais situações o ORA
ainda é preferível? (c) Considere um experimento com apenas 3 amostras por
condição --- o GSEA ainda é confiável? Por quê?
\end{exercicio}

\begin{exercicio}
\textbf{Modelo da trealose.} No modelo de quatro EDOs descrito no capítulo
para o ciclo da trealose: (a) Analise a estabilidade do estado estacionário
não-estressado ($[\text{Trealose}] = 5\;\mathrm{mM}$). (b) Submeta o sistema
a um pulso de estresse e simule durante 2 h. Compare com os dados
experimentais da tabela. (c) Plote o gráfico resultante e calcule o RMSE
do ajuste. (d) Repita o ajuste com barreiras de 50\% nos parâmetros e
discuta a identificabilidade do modelo.
\end{exercicio}

\begin{exercicio}
\textbf{Parâmetros distribuídos.} Modifique o modelo da trealose para
permitir variabilidade intra-populacional: sortear 100 conjuntos de
parâmetros a partir de distribuições gaussianas em torno dos valores
estimados (10\% de variância). (a) Plote a distribuição do tempo de meia
para atingir 50 mM de trealose. (b) Quantas células atingem o estado
estacionário pré-estresse no tempo de 120 min? (c) Discuta implicações para
estudos single-cell.
\end{exercicio}

\begin{exercicio}
\textbf{FBA com expressão.} Carregue o modelo iMM904 de \textit{S.\
cerevisiae} usando COBRApy. (a) Calcule a taxa de crescimento ótima por FBA
sem restrições de expressão. (b) Implemente o método E-Flux descrito no
texto e recalcule a taxa de crescimento. (c) Identifique as 10 reações
mais reprimidas por baixa expressão. (d) Compare a biomassa e a produção
de trealose em ambas as configurações. (e) Valide com dados de
literatura sobre crescimento de \textit{S. cerevisiae} sob estresse
térmico.
\end{exercicio}

\begin{exercicio}
\textbf{Análise de redes.} Usando uma matriz de correlação de Pearson com
1000 genes de \textit{S. cerevisiae} (gerada por simulação ou obtida do
GEO): (a) Construa uma rede de co-expressão com $|r| > 0{,}8$ e
$p < 0{,}01$ (correção de Bonferroni). (b) Identifique módulos pelo
algoritmo Louvain usando \texttt{NetworkX}. (c) Calcule as centralidades de
grau, betweenness e eigenvector para todos os nós. (d) Identifique os 20
hubs por centralidade de grau e analise enriquecimento GO para esse
conjunto. (e) Comente sobreposição com fatores de transcrição conhecidos
em levedura.
\end{exercicio}

\begin{exercicio}
\textbf{Modelo híbrido.} Construa um modelo híbrido para \textit{S.\
cerevisiae} combinando: (i) regulação transcricional de três genes do
ciclo da trealose (\textit{TPS1}, \textit{TPS2}, \textit{NTH1}) por ODEs;
(ii) metabolismo central por FBA do modelo iMM904; (iii) acoplamento
via restrições de expressão (modificação de bounds da reação TPS1_rxn
proporcional a $[\textit{TPS1}]$). (a) Escreva o sistema de EDOs acoplado.
(b) Implemente em Python. (c) Simule um pulso de choque térmico a 37°C
durante 2 h e meça o crescimento e a produção de trealose. (d) Identifique
predições não-triviais do modelo (acoplamento metabolização-tempo de
formação de trealose).
\end{exercicio}

\begin{exercicio}
\textbf{Projeto integrador.} Escolha um sistema biológico diferente do
capítulo --- por exemplo, o sistema de defesa antioxidante em humanos ou a
via de biossíntese de flavonoides em plantas --- e: (a) Formule um conjunto
de hipóteses biológicas integrando o que se sabe por cada camadaômica;
(b) projete um experimento multi-ômico apropriado; (c) construa um modelo
matemático com 3--5 variáveis; (d) implemente o modelo em Python;
(e) analise a sensibilidade dos parâmetros; (f) formule previsões falsificáveis
sobre mutantes e drogas; (g) proponha uma estratégia de validação
experimental completa.
\end{exercicio}

\section{Notas Bibliográficas}
\label{sec:11-bibliografia}

{A integração multi-ômica em biologia de sistemas tem três textos-fonte
consolidados. \textit{Systems Biology: A Textbook} de Klipp, Herrmann,
Liebermeister, Rother, Gilles, Wierling, Kowald e Lehrach (2ª edição,
Wiley-VCH) é o tratamento mais pedagógico do ferramental --- cobre
estequiometria, termodinâmica, cinética e integração multiômica com
exemplos detalhados. \textit{Systems Biology: Constraint-based
Reconstruction and Analysis} de Bernhard Palsson (Cambridge) é a referência
para a abordagem GEM/FBA: estequiometria, espaços de soluções, integração
com dados transcriptômicos (E-Flux, GIMME, iMAT). O artigo seminal de
Ideker, Galitski e Hood (\textit{A new approach to decoding life: systems
biology}, Annual Review of Genomics and Human Genetics, 2001) formulou
a agenda do campo --- identificar módulos responsivos como unidades de
organização biológica --- e permanece leitura introdutória recomendada.

Para a frente computacional e clínica, o artigo de revisão de Hasin,
Seldin e Lusis (\textit{Multi-omics approaches to disease}, Genome Biology,
2017) sumariza o estado da arte até aquela data e oferece uma taxonomia
útil para integração em doenças. Ritchie et al.\ (Nature Reviews Genetics,
2015) discutem métodos de integração genética-ômica, com ênfase em modelos
multi-ômicos de fenótipos complexos. Huang, Chaudhary e Lee (\textit{More is
better: recent progress in multi-omics data integration methods}, Frontiers
in Genetics, 2017) oferecem um panorama focando em MOFA e ferramentas
relacionadas. Para enriquecimento de vias, o artigo clássico de
Subramanian et al.\ (2005) introduzindo o GSEA é leitura obrigatória;
GSEApy é a implementação Python recomendada. Para o cenário específico
de integração em escala atlas --- quando se deseja comparar métodos
de scRNA-seq ou scATAC-seq sobre dados com múltiplos lotes, laboratórios
e condições ---, o benchmark de Luecken e
colaboradores~\cite{luecken2022} oferece um protocolo padronizado
(catorze métricas, três dimensões de avaliação) e uma recomendação
operacional dos métodos com melhor desempenho em tarefas complexas.

Os bancos de dados e ferramentas computacionais evoluem rapidamente.
COBRApy (\url{https://opencobra.github.io/cobrapy/}) é o pacote padrão de
Python para análise metabólica em escala genômica; NetworkX
(\url{https://networkx.org}) é referência para análise de redes; Enrichr
(\url{https://maayanlab.cloud/Enrichr/}) é uma plataforma web de análise
de enriquecimento cobrindo uma centena de bibliotecas gênicas; STRING
oferece redes de interação proteína-proteína curadas e complementares às
redes inferidas dos dados; KEGG e Reactome são as duas referências de via
mais usadas. Cursos de pós-graduação em bioinformática e biologia
computacional --- MIT, EMBL-EBI, Coursera --- oferecem treinamentos
regulares sobre essas ferramentas.

Os livros complementares --- \textit{An Introduction to Systems Biology} de
Uri Alon, \textit{A First Course in Systems Biology} de Eberhard Voit e
\textit{Fundamentals of Systems Biology} de Markus Covert --- oferecem
perspectivas pedagógicas distintas. Alon enfatiza os princípios de design;
Voit detalha a formulação bioquímica dos modelos; Covert cobre a
integração experimental e computacional.

}

\begin{thebibliography}{99}
\bibitem{klipp} Klipp E., Herrmann H., Liebermeister W., Rother K., Gilles
E.D., Wierling C., Kowald A., Lehrach H. (2016). \textit{Systems Biology:
A Textbook}, 2ª ed. Wiley-VCH, Weinheim.

\bibitem{palsson} Palsson B.O. (2015). \textit{Systems Biology:
Constraint-based Reconstruction and Analysis}. Cambridge University Press,
Cambridge.

\bibitem{ideker} Ideker T., Galitski T., Hood L. (2001). A new approach to
decoding life: systems biology. \textit{Annual Review of Genomics and Human
Genetics}, 2:343--372.

\bibitem{hasin} Hasin Y., Seldin M., Lusis A. (2017). Multi-omics approaches
to disease. \textit{Genome Biology}, 18:83.

\bibitem{subramanian} Subramanian A., Tamayo P., Mootha V.K., Mukherjee S.,
Ebert B.L., Gillette M.A., Paulovich A., Pomeroy S.L., Golub T.R., Lander
E.S., Mesirov J.P. (2005). Gene set enrichment analysis: a
knowledge-based approach for interpreting genome-wide expression profiles.
\textit{Proceedings of the National Academy of Sciences USA}, 102:15545--15550.

\bibitem{ritchie} Ritchie M.D., Holzinger E.R., Li R., Pendergrass S.A.,
Kim D. (2015). Methods of integrating data to uncover genotype--phenotype
interactions. \textit{Nature Reviews Genetics}, 16:85--97.

\bibitem{huang} Huang S., Chaudhary K., Lee J.K. (2017). More is better:
recent progress in multi-omics data integration methods. \textit{Frontiers
in Genetics}, 8:84.

\bibitem{alon} Alon U. (2019). \textit{An Introduction to Systems Biology:
Design Principles of Biological Circuits}, 2ª ed. CRC Press, Boca Raton.

\bibitem{voit} Voit E.O. (2017). \textit{A First Course in Systems Biology},
2ª ed. Garland Science, Nova York.

\bibitem{covert} Covert M.W. (2014). \textit{Fundamentals of Systems
Biology}. CRC Press, Boca Raton.

\bibitem{cavill} Cavill R., Jennen D., Kleinjans J., Briede J.J. (2016).
Transcriptomic and metabolomic data integration. \textit{Briefings in
Bioinformatics}, 17:891--901.

\bibitem{pronk} Dijkgraaf J., Brown R., O'Neil J., Van Dijken J., Pronk J.T.,
Geertsma Y.R. (2014). Metabolic engineering for trehalose accumulation in
yeast. \textit{Applied and Environmental Microbiology}, 80:143--150.

\bibitem{dinh} Dinh H.V., Niu W., Chen Y., Smith R.P., Ko A., Lu W., Chen K.,
Sander C., Price N.D., Palsson B.O. (2022). An integrated multi-scale
model of the yeast metabolic regulation network. \textit{Molecular Systems
Biology}, 18:e10806.

\bibitem{cobra} Ebrahim A., Lerman J.A., Palsson B.O., Hyduke D.R. (2013).
COBRApy: constraints-based reconstruction and analysis for python. \textit{BMC
Systems Biology}, 7:74. Disponível em \url{https://opencobra.github.io/cobrapy/}.

\bibitem{enrichr} Kuleshov M.V., Jones M.R., Rouillard A.D., Fernandez N.F.,
Duan Q., Wang Z., Koplev S., Jenkins S.L., Jagannathan K., Lachmann A.,
McDermott M.G., Monteiro C.D., Gundersen G.W., Ma'ayan A. (2016). Enrichr:
a comprehensive gene set enrichment analysis web server 2016 update.
\textit{Nucleic Acids Research}, 44:W90--W97. Disponível em
\url{https://maayanlab.cloud/Enrichr/}.

\bibitem{gseapy} Subramanian A., Kuehn H., Gould J., Tamayo P., Mesirov J.P.
(2007). GSEA-P: a desktop application for Gene Set Enrichment Analysis.
\textit{Bioinformatics}, 23:3251. Disponível em
\url{https://gseapy.readthedocs.io}.

\bibitem{luecken2022} Luecken M.D., Büttner M., Chaichoompu K., Danese A.,
Interlandi M., Mueller M.F., Strobl D.C., Zappia L., Dugas M.,
Colomé-Tatché M., Theis F.J. (2022). Benchmarking atlas-level data
integration in single-cell genomics. \textit{Nature Methods},
19:41--50. Disponível em \url{https://doi.org/10.1038/s41592-021-01336-8}.
\end{thebibliography}
